\section{Stato dell'arte}

\subsection{Dal contesto conversazionale alla query di ricerca}
Nel recupero conversazionale, la pertinenza di un passaggio dipende dal significato che il turno corrente assume nel dialogo. QReCC studia questo problema associando alle domande contestuali riscritture autonome, utili per valutare separatamente interpretazione della domanda, recupero e risposta \cite{qrecc}. RECOR conserva la dipendenza tra turni e aggiunge casi in cui, anche una volta risolto il riferimento, la relazione tra domanda e passaggio richiede un collegamento concettuale \cite{recor}. Distinguere i due passaggi è importante: esplicitare a quale entità si riferisce un pronome può facilitare la ricerca, mentre capire perché un documento su un'entità diversa risponda alla domanda può richiedere informazioni che emergono soltanto dal confronto con quel documento.

La rappresentazione più semplice accoda la storia alla domanda e lascia al retriever il compito di estrarne il segnale pertinente. È un controllo sostanziale: nelle baseline ufficiali di RETECO, sullo stesso dev e con BM25, aggiungere la storia aumenta la nDCG@10 macro da 0,1827 a 0,4379 \cite{data}. La concatenazione, tuttavia, può conservare riferimenti a temi abbandonati o competere con il turno corrente per il budget di token. ConvGQR combina una riscrittura con una risposta ipotetica e introduce un segnale legato al recupero per orientare la riformulazione \cite{convgqr}. Una risposta ipotetica può suggerire termini assenti dalla storia, ma può anche introdurre ipotesi errate che spostano il ranking. ConvSearch-R1 ottimizza la riscrittura mediante apprendimento per rinforzo e usa la posizione dei documenti pertinenti come segnale di ricompensa \cite{convsearchr1}. Questo approccio lega la riformulazione al retriever usato per ottenere la ricompensa. La nostra variante riscritta usa un modello istruito, con la formulazione contestuale come controllo. Il miglioramento riportato da RECOR con storia e informazioni esplicite di ragionamento motiva il confronto \cite{recor}. L'effetto di una riscrittura prodotta automaticamente va misurato per ciascun retriever.

Una seconda famiglia apprende direttamente una rappresentazione densa del turno e della sua storia. ConvDR distilla nel codificatore conversazionale le rappresentazioni prodotte da un insegnante che osserva riscritture di riferimento, riutilizzando il codificatore dei documenti \cite{convdr}. HAConvDR seleziona invece informazioni utili dalla cronologia e ricava segnali di supervisione dall'effetto dei turni precedenti sul recupero, affrontando in modo esplicito le conversazioni lunghe con cambi di argomento \cite{haconvdr}. I due lavori offrono un riferimento per distinguere tre meccanismi: migliorare il testo da cercare, selezionare il contesto pertinente e apprendere come rappresentarlo. Nel nostro bi-encoder la supervisione riguarda la pertinenza dei documenti, attraverso esempi positivi e negativi.

\subsection{Generazione dei candidati e riordinamento}
BM25 valorizza le corrispondenze lessicali e può essere efficace quando nomi, acronimi o termini specialistici identificano il passaggio pertinente \cite{bm25}. Un bi-encoder rappresenta domanda e documento separatamente e può recuperare passaggi con formulazioni diverse. Il confronto tra vettori precalcolati consente di cercare su tutto il corpus, a fronte del costo iniziale di codifica. Quando si aggiorna il codificatore dei documenti, occorre ricostruire le relative rappresentazioni. Le differenze fra i due segnali giustificano l'uso di Reciprocal Rank Fusion (RRF): la fusione combina le posizioni senza richiedere punteggi calibrati sulla stessa scala \cite{rrf}. La sua efficacia dipende da quali passaggi i due retriever portano in graduatoria e con quale posizione.

Una volta fissati i candidati, il riordinatore può valutare congiuntamente domanda e testo. Qwen3-Reranker assegna un punteggio a ogni coppia mediante il confronto tra i logit di due token di giudizio \cite{qwen}. Un approccio a lista come RankGPT chiede invece a un modello generativo di produrre una permutazione dei candidati \cite{rankgpt}. Quest'ultimo può confrontare più passaggi nello stesso prompt, mentre la finestra di contesto limita quanti ne può esaminare insieme e la generazione aggiunge costo. Il riordinamento per coppia rende invece naturale ripetere il punteggio sugli stessi input, con lo stesso budget di candidati. Per il 2a questo permette di attribuire una variazione di nDCG@10 alla nuova disposizione dei passaggi. Il Recall@$K$ del pool descrive quanta evidenza pertinente è accessibile al riordinatore prima dell'intervento.

\subsection{Generazione fondata su passaggi}
Nei sistemi di domanda e risposta aperta, Fusion-in-Decoder mostra come un modello sequenza-a-sequenza possa aggregare evidenze codificate separatamente per ciascun passaggio \cite{fid}. ALCE affronta un aspetto complementare: valuta in che misura le citazioni generate sostengano effettivamente le affermazioni della risposta \cite{alce}. Nel 2b di RETECO l'uso di passaggi forniti dagli organizzatori sposta il confronto sulla capacità del generatore di interpretare, selezionare e combinare il materiale disponibile. La rubrica distingue correttezza, completezza, pertinenza, coerenza conversazionale e fedeltà \cite{evaluation}. Per esempio, una risposta può riportare un dato corretto in termini generali ma assente dai passaggi assegnati. Può anche limitarsi a informazioni ben supportate lasciando irrisolta una parte della domanda.

Il nostro generatore usa un prompt con storia, turno corrente ed evidenze. L'adattamento tramite LoRA e quantizzazione a 4 bit interviene sul costo della specializzazione del modello \cite{lora,qlora}. La ripetizione dello strato interviene invece sul costo e sul comportamento dell'inferenza. Un confronto fra modello base e adattato a profondità originale misura il contributo dell'adattamento. Ripetere il confronto di profondità sul medesimo checkpoint adattato, con gli stessi passaggi e la stessa decodifica, permette di studiare il secondo effetto.

\subsection{Profondità ricorrente e posizione del presente studio}
Negli Universal Transformers una trasformazione condivisa aggiorna iterativamente le rappresentazioni lungo la profondità \cite{universal}. Geiping e collaboratori addestrano un modello linguistico con un blocco ricorrente il cui numero di applicazioni può essere aumentato al momento dell'inferenza \cite{geiping}. In entrambi i casi, la ricorrenza entra nella progettazione e nell'addestramento. I risultati di queste architetture motivano lo studio del calcolo nello spazio latente, ma il comportamento di un blocco addestrato a iterare non determina quello di uno strato standard riutilizzato dopo l'addestramento.

I precedenti più diretti sono gli interventi sui checkpoint esistenti. Lys e collaboratori riapplicano in inferenza un intervallo di strati e analizzano le traiettorie delle rappresentazioni intermedie \cite{lys}. Chen e collaboratori mostrano che riapplicare direttamente un blocco può peggiorare i risultati e propongono di frazionare le variazioni residue in aggiornamenti smorzati \cite{chen}. Nel nostro esperimento la riapplicazione diretta di un singolo strato centrale offre il controllo più semplice: si varia il numero di applicazioni con pesi, input e posizione dello strato fissati. Questa restrizione rende verificabile il contrasto tra una e due applicazioni anche quando il miglioramento atteso è piccolo o quando il risultato è negativo.

La Tabella~\ref{tab:related_techniques} colloca l'intervento rispetto alle principali tecniche. Il contributo del presente studio consiste nel valutare la stessa modifica di profondità in due punti della pipeline, sul giudizio di pertinenza nel 2a e sulla risposta autoregressiva nel 2b. Nel primo caso il costo cresce con il numero di coppie valutate. Nel secondo intervengono anche la lunghezza del contesto e il numero di token generati. Per questo qualità e latenza vanno lette insieme, usando candidati o passaggi fissati e registrando le condizioni di inferenza descritte nella metodologia.

\begin{table}[!htbp]
\centering
\small
\caption{Confronto metodologico tra tecniche per il Track 2 di RETECO. I costi indicano le operazioni aggiuntive introdotte da ciascun meccanismo.}
\label{tab:related_techniques}
\setlength{\tabcolsep}{4pt}
\renewcommand{\arraystretch}{1.24}
\begin{tabular}{@{}>{\raggedright\arraybackslash}p{.15\linewidth}>{\raggedright\arraybackslash}p{.21\linewidth}*{2}{>{\raggedright\arraybackslash}p{.28\linewidth}}@{}}
\toprule
\textbf{Stadio} & \textbf{Tecnica} & \textbf{Operazione} & \textbf{Costo o limite} \\
\midrule
Query (2a) & Storia con BM25 \cite{bm25,data} & Cerca il turno insieme alla cronologia in un indice invertito. & Termini dei turni precedenti possono introdurre corrispondenze estranee. \\
Query (2a) & Riscrittura generativa \cite{qrecc,convgqr,convsearchr1} & Produce una query autonoma e può espandere il vocabolario della ricerca. & Richiede una generazione per turno. Errori nella riscrittura si propagano ai candidati. \\
Candidati (2a) & Bi-encoder conversazionale \cite{convdr,haconvdr,qwen} & Confronta la rappresentazione della query con vettori dei passaggi precalcolati. & L'adattamento può richiedere esempi di pertinenza. Aggiornare il codificatore dei passaggi impone una nuova codifica del corpus. \\
Candidati (2a) & Fusione RRF \cite{rrf} & Aggrega ranking lessicali e densi usando le posizioni. & Richiede entrambe le ricerche. I documenti assenti dalle due liste rimangono esclusi. \\
Ranking (2a) & Punteggio per coppia \cite{qwen} & Valuta insieme la query e ciascun passaggio candidato. & Il numero di forward dipende dal numero di candidati e dal batch. \\
Ranking (2a) & Generazione di una lista \cite{rankgpt} & Ordina più candidati nello stesso prompt producendo una permutazione. & Il contesto limita la dimensione della lista per prompt. Serve decodificare l'ordinamento. \\
Risposta (2b) & Generazione con evidenze \cite{fid,alce,evaluation} & Combina passaggi di supporto per rispondere al turno. & Contano il budget dei passaggi e la lunghezza della risposta. La fedeltà richiede una misura dedicata. \\
Decoder & Ricorrenza appresa \cite{universal,geiping} & Addestra un modello che applica più volte una trasformazione condivisa. & Richiede addestramento con la ricorrenza e più calcolo quando si aumenta la profondità. \\
Decoder & Ricorrenza a inferenza \cite{lys,chen} & Riutilizza strati di un checkpoint esistente, con eventuale smorzamento. & Aggiunge applicazioni dello strato. L'effetto dipende dal blocco e dalla regola di aggiornamento. \\
Decoder (2a/2b) & Ripetizione diretta di uno strato (questo studio) & Varia la profondità del riordinatore o del generatore con gli stessi input. & Con due passaggi aggiunge un'applicazione per forward. Il confronto generativo usa la cache chiave--valore disabilitata. \\
\bottomrule
\end{tabular}
\end{table}
